% ! TEX root = ../am2-19-20.tex

\chapter{Numere și funcții complexe --- recapitulare}

\section{Numere complexe -- Noțiuni de bază}

Începem cu cîteva noțiuni esențiale și recapitulative privitoare la mulțimea
numerelor complexe. Amintim definiția:
\[
  \CC = \{ a + bi \mid a, b \in \RR, i^2 = -1 \},
\]
precum și faptul că avem, în general, șirul de incluziuni:
\[
  \NN \seq \ZZ \seq \QQ \seq \RR \seq \CC.
\]

Dat un număr complex $ z = a + bi $, se numește \emph{partea reală},
notată $ \dr{Re}(z) $, numărul real $ a $, iar \emph{partea imaginară},
notată $ \dr{Im}(z) $, numărul real $ b $.

De asemenea, \emph{conjugatul} numărului complex $ z $ de mai sus este
$ \overline{z} = z^\ast = a - bi $, iar \emph{modulul} numărului complex
$ z $ este $ |z| = \sqrt{a^2 + b^2} $.

Există mai multe forme de reprezentare a numerelor complexe:
\begin{itemize}
\item \emph{forma algebrică}, dată mai sus, $ z = a + bi, a, b \in \RR $;
\item \emph{forma trigonometrică}, $ z = r(\cos \theta + i \sin \theta) $,
  unde $ r = |z| $, iar $ \theta = \dr{Arg}(z) $ se numește \emph{argumentul %
  principal};
\item \emph{forma polară}, $ z = r e^{i\theta} = r \exp(i \theta) $, unde
  $ r $ și $ \theta $ au înțelesul din forma trigonometrică;
\item \emph{forma geometrică}, în care $ z = a + bi $ reprezintă
  \emph{afixul} punctului din plan $ A(a, b) $. Pentru acest caz, menționăm
  că $ |z| = \sqrt{a^2 + b^2} $ reprezintă lungimea vectorului de poziție
  al punctului $ A $, iar $ \theta $ reprezintă unghiul pe care îl face
  acest vector de poziție cu axa $ OX $, măsurat în sens trigonometric.
\end{itemize}

Operațiile cu numere complexe se fac în modul uzual, ținînd seama de
proprietatea $ i^2 = -1 $. Mai amintim \emph{formula lui Moivre}, utilă
în special atunci cînd numărul complex a fost scris sub formă trigonometrică.
Fie, așadar, numerele complexe:
\[
  z_1 = r_1(\cos \theta_1 + i\sin\theta_1), \quad %
  z_2 = r_2(\cos \theta_2 + i \sin\theta_2).
\]
Atunci înmulțirea acestora se face cu formula:
\[
  z_1 \cdot z_2 = r_1r_2 \left( \cos(\theta_1 + \theta_2) + i \sin(\theta_1 + \theta_2) \right),
\]
formulă care se generalizează ușor în forma:
\[
  z^n = r^n(\cos (n\theta) + i \sin (n\theta)), \forall n.
\]

Tot folosind numere complexe, putem reprezenta și curbe:
\begin{itemize}
\item \emph{Cercul} centrat în punctul de afix $ z_0 $ și de rază $ R $ are ecuația:
  \[
    |z - z_0| = r;
  \]
\item \emph{Elipsa} cu focarele în punctele de afixe $ z_0 $ și $ w_0 $,
  iar axa mare are lungimea $ d $ are ecuația:
  \[
    |z - z_0| + |z - w_0| = d.
  \]
\end{itemize}

Amintim și \emph{formele canonice} ale ecuațiilor acestor conice:
\begin{itemize}
\item \emph{Cercul} centrat în $ (x_0, y_0) $ și de rază $ R $ are
  ecuația:
  \[
    (x - x_0)^2 + (y - y_0)^2 = R^2;
  \]
\item \emph{Elipsa} de semiaxe $ a, b $ are ecuația:
  \[
    \dfrac{x^2}{a^2} + \dfrac{y^2}{b^2} = 1.
  \]
\end{itemize}

Tot din punct de vedere geometric, mai amintim și că \emph{distanța} între
două puncte $ A(z) $ și $ B(w) $ este $ AB = |z - w| $.

%%%%%%%%%%%%%%%%%%%%%%%%%%%%%%%%%%%%%%%%%%%%%%%%%%%%%%%%%%%%%%%%%%%%%%
\section{Funcții complexe elementare}

Cea mai simplă funcție complexă este \emph{funcția exponențială}, definită prin:
\[
  \exp : \CC \to \CC, \quad \exp z = e^z = \sum_{n \geq 0} \dfrac{z^n}{n!}.
\]

Se observă imediat că au loc proprietățile așteptate:
\begin{itemize}
\item $ \exp(0) = 1 $;
\item $ \exp(z_1 + z_2) = \exp(z_1) \cdot \exp(z_2), \forall z_{1,2} \in \CC $;
\item $ \exp(iy) = \cos y + i \sin y, \forall y \in \RR $ (formula lui Euler).
\end{itemize}

Mai departe, putem calcula simplu \emph{logaritmul complex}, dacă numărul complex
a fost adus în forma polară. Fie, așadar, $ z = re^{i\theta} $. Rezultă:
\[
  \ln z = \ln r + i\theta.
\]

În general, cum argumentul unui număr complex nu este unic ($ \theta $ de mai sus
reprezintă \emph{argumentul principal}, dar $ \theta + 2 k\pi $ este argumentul
general), spunem că funcția logaritm este \emph{multiformă}, deoarece valoarea ei
generală este:
\[
  \dr{Ln}z = \{ \ln |z| + i(\dr{Arg}z + 2 k \pi) \mid k \in \ZZ \},
\]
iar $ \ln z $ se numește \emph{valoarea principală} a logaritmului.

\emph{Funcția putere} se definește acum simplu pornind de la formula:
\[
  a^b = \exp\left(b \ln a\right).
\]

Rezultă:
\[
  z^m = \exp\left(m \ln z\right) = \exp \left( m(\ln|z| + i\dr{Arg}z) \right),
\]
folosind valoarea principală.

Similar se definește și puterea rațională, adică \emph{funcția radical}:
\[
  \sqrt[n]{z} = z^{\frac{1}{n}} = \exp\left( \dfrac{1}{n} \ln z \right).
\]

Folosind funcțiile exponențiale și identitatea lui Euler, putem defini și
\emph{funcții trigonometrice complexe:}
\begin{align*}
  \cos z &= \dfrac{\exp(iz) + \exp(-iz)}{2} \\
  \sin z &= \dfrac{\exp(iz) - \exp(-iz)}{2i} \\
  \tan z &= -i\dfrac{\exp(iz) - \exp(-iz)}{\exp(iz) + \exp(-iz)} \\
\end{align*}

Mai avem nevoie și de \emph{funcțiile trigonometrice hiperbolice}:
\begin{align*}
  \sinh z &= \dfrac{\exp(z) - \exp(-z)}{2} \\
  \cosh z &= \dfrac{\exp(z) + \exp(-z)}{2} \\
  \tanh z &= \dfrac{\sinh z}{\cosh z}
\end{align*}
și au loc legăturile:
\[
  \sinh z = -i\sin(iz), \quad \cosh z = \cos(iz).
\]

\emph{Funcțiile trigonometrice inverse} se pot obține din rezolvarea unor
ecuații trigonometrice\footnotemark:
\[
  w = \dr{arcsin z} \Rightarrow z = \sin w \Rightarrow %
  z = \dfrac{\exp(iw) - \exp(-iw)}{2i} \Leftrightarrow %
  \exp(2iw) - 2iz\exp(iw) - 1 = 0,
\]
care se rezolvă (pentru $ w $) ca o ecuație de gradul al doilea și se obține:
\[
  w = \dr{arcsin} z = -i\ln(iz \pm \sqrt{1 - z^2})
\]
și se procedează similar pentru $ \dr{arccos} $ și $ \dr{arctan} $.

\footnotetext{Pentru corectitudine, ar trebui să notăm funcțiile
  trigonometrice cu inițială mare, $ \dr{Arcsin}, \dr{Arccos} $ etc. Dar
  vom folosi de fiecare dată doar valoarea principală, astfel că, prin
  abuz de notație, folosim scrierea din cazul real. Însă trebuie reținut
  că majoritatea funcțiilor complexe sînt \emph{multiforme}, i.e.\ pot
  avea mai multe valori!}

%%%%%%%%%%%%%%%%%%%%%%%%%%%%%%%%%%%%%%%%%%%%%%%%%%%%%%%%%%%%%%%%%%%%%%

\section{Funcții olomorfe}

Fie o funcție complexă oarecare $ f : A \to \CC $, cu $ A \seq \CC $.
Pentru orice $ z \in A \seq \CC $, deoarece $ z $ are o parte reală și
o parte imaginară, și imaginea sa prin $ f $ se poate separa. Deci, în general,
orice funcție complexă $ f $ ca mai sus poate fi scrisă sub forma
\[
  f = P + iQ, \quad P, Q : \RR^2 \to \RR.
\]

Rezultă că noțiunile de \emph{limită, continuitate, derivabilitate} pot fi puse
pe componente.

\begin{definition}\label{def:olo}
  O funcție complexă $ f : \CC \to \CC $ se numește \emph{olomorfă} dacă este
  derivabilă în orice punct din domeniul de definiție.
\end{definition}

Nu intrăm în detalii, deoarece nu vom rezolva exerciții cu derivate complexe.

Vor fi foarte importante, însă, rezultatele:

\begin{theorem}\label{thm:cauchy-riemann}
  Funcție $ f : \CC \to \CC, f = P + iQ $ este olomorfă dacă $ P, Q : \RR^2 \to \RR $
  sunt diferențiabile, iar derivatele lor parțiale verifică \emph{condițiile %
    Cauchy-Riemann}, adică:
  \[
    \dfrac{\partial P}{\partial x} = \dfrac{\partial Q}{\partial y}, \quad %
    \dfrac{\partial Q}{\partial x} = -\dfrac{\partial P}{\partial y}.
  \]
\end{theorem}

\begin{remark}\label{rk:der-part}
  Pentru simplitate, vom mai nota derivatele parțiale cu indici, adică, de exemplu:
  \[
    \dfrac{\partial f}{\partial x} \stackrel{\dr{not.}}{=\joinrel=} f_x
  \]
  și similar $ f_y, f_{xx}, f_{xy}, f_{yx} $ etc.
\end{remark}

\begin{corollary}\label{cor:olo-arm}
  Dacă funcția complexă $ f = P + iQ $ este olomorfă, atunci $ P $ și $ Q $ sînt
  armonice, adică:
  \[
    \Delta P = P_{xx} + P_{yy} = \Delta Q = Q_{xx} + Q_{yy} = 0.
  \]
\end{corollary}

Atenție, însă, la formularea rezultatului: condiția de olomorfie este \emph{necesară},
în niciun caz suficientă! Negarea corolarului de mai sus este:

\begin{corollary}\label{cor:not-olo}
  Dacă una dintre funcțiile $ P $ sau $ Q $ nu este armonică, atunci funcția
  $ f = P + iQ $ nu poate fi olomorfă.
\end{corollary}

%%%%%%%%%%%%%%%%%%%%%%%%%%%%%%%%%%%%%%%%%%%%%%%%%%%%%%%%%%%%%%%%%%%%%%

\section{Exerciții}

\indent\indent 1. Verificați dacă funcția de mai jos este olomorfă:
\[
  f : \CC \to \CC, \quad f(z) = z^2 + \exp(iz).
\]

\emph{Indicație:} Se separă partea reală și partea imaginară a funcției
și se verifică condițiile Cauchy-Riemann și faptul că cele două componente
sînt armonice.

\vspace{0.3cm}

2. Fie $ P(x, y) = e^{2x}\cos 2y + y^2 - x^2 $. Determinați funcția olomorfă
$ f = P + iQ $ astfel încît $ f(0) = 1 $.

\emph{Indicație:} Verificăm dacă $ \Delta P = 0 $ (condiție necesară!).
Apoi, prin integrarea condițiilor Cauchy-Riemann, se obține componenta
$ Q $. În final, $ f(z) = e^{2z} - z^2 + ki, k \in \CC $ și folosind
condiția din enunț, obținem $ k = 0 $.

\vspace{0.3cm}

3. Rezolvați ecuațiile:
\begin{enumerate}[(a)]
\item $ \exp w = -2i $;
\item $ z^3 + 2 - 2i = 0 $;
\item $ \sin z = 2 $.
\end{enumerate}

\vspace{0.3cm}

4. Calculați:
\begin{enumerate}[(a)]
\item $ \sin(1 + i) $;
\item $ \sinh(1 - i) $;
\item $ \ln i $;
\item $ \ln(1 - i) $;
\item $ (1 + i)^{20} $;
\item $ \sqrt[5]{1 - i} $;
\item $ \dr{arcsin}(i\sqrt{3}) $;
\item $ \dr{arccos}(i\sqrt{3}) $;
\item $ \tan(1 - i) $.
\end{enumerate}

\vspace{0.3cm}

5. Determinați funcția olomorfă $ f(z) = P(x, y) + iQ(x, y) $, pentru:
\begin{enumerate}[(a)]
\item $ P(x, y) = x^2 - y^2 - 2y $, știind că $ f(0) = i $;
\item $ P(x, y) = x^4 - 6x^2y^2 + y^4 $, știind că $ f(1) = 1 $;
\item $ P(x, y) = (x \cos y - y \sin y)e^x $.
\end{enumerate}

\vspace{0.3cm}

6. Scrieți sub formă trigonometrică și polară numerele complexe și
reprezentați-le grafic:
\begin{enumerate}[(a)]
\item $ z = 3 - i $;
\item $ z = 3 + i $;
\item $ z = i $;
\item $ z = 1 $;
\item $ z = 1 + 2i $;
\item $ z = 2 + i $.
\end{enumerate}

\vspace{0.3cm}

7. Găsiți forma canonică și ecuația complexă pentru:
\begin{enumerate}[(a)]
\item Cercul centrat în $ (0, 1) $ și cu raza $ 2 $;
\item Cercul centrat în $ (1, 0) $ și cu raza $ 1 $;
\item Cercul centrat în $ (1, 2) $ și cu raza $ 1 $;
\item Elipsa cu focarele în $ (-1, 0) $ și $ (3, 0) $ și cu
  axa mare de lungime 6;
\item Elipsa cu focarele în $ (0, 1) $ și $ (0, -2) $ și
  cu axa mare de lungime 5.
\end{enumerate}

Reprezentați grafic fiecare dintre cazurile de mai sus.

\vspace{0.3cm}

8. Fie punctele $ A(1 + 2i) $ și $ B(-1) $, iar $ M $, mijlocul
segmentului $ [AB] $. Calculați distanța de la punctul $ M $
la punctul $ N $, de afix $ -2 + 3i $. Reprezentare grafică.

\vspace{0.3cm}

\textbf{Observație:} De data aceasta, pentru temă NU mai este suficient să
rezolvați un singur subpunct de la un exercițiu, cel puțin nu pentru toate
exercițiile. Astfel, tema poate conține oricare dintre variantele de mai jos:
\begin{itemize}
\item exercițiul 1;
\item exercițiul 2;
\item un subpunct de la exercițiul 3 și unul de la exercițiul 4;
\item două subpuncte care folosesc funcții diferite de la exercițiul 4
  (exemple: c și f, g și e, a și f etc.);
\item un subpunct de la exercițiul 5;
\item un subpunct de la exercițiul 6 și un subpunct de la exercițiul 7;
\item un subpunct de la exercițiul 7 și exercițiul 8.
\end{itemize}

\newpage

\textbf{Temă specială:} Propun și o temă specială pentru această lecție.
Alegeți un subiect din cartea lui Paul Nahin despre istoria numerelor complexe
și faceți un clip în care să îl prezentați. Detalii și cerințe:
\begin{itemize}
\item puteți alege orice subiect din carte, DAR trebuie să mă anunțați cînd ați
  ales prin email și să-l validez eu;
\item după validare, aveți maximum 2 săptămîni la dispoziție;
\item videoclipul trebuie să dureze minimum 10 minute;
\item conținutul matematic sofisticat nu este obligatoriu; puteți insista pe aspecte
  istorice;
\item puteți folosi și alte surse, dar subiectul să fie legat de istoria descoperirilor
  și aplicațiilor numerelor complexe;
\item dacă nu știți ce subiect să alegeți, scrieți-mi și vi-l propun eu.
\end{itemize}

\textbf{Punctaje:
\begin{itemize}
\item video $ \geq $ 10 minute = 5 puncte de seminar și $ \geq $ 5, $ \leq $ 10 puncte la examenul
  final;
\item eseu $ \geq $ 5 pagini = 5 puncte de seminar;
\item video + eseu = 8 puncte la seminar și $ \geq $ 5, $ \leq $ 10 puncte la examenul final\footnotemark.
\end{itemize}
}

\footnotetext{Asta înseamnă că \emph{sigur} rotunjim în sus de la o notă $ x,5 $ și,
  în funcție de cum vă descurcați la examen, \emph{este posibil} să rotunjim chiar și de
  la $ x,2 $ sau $ x,3 $ sau $ x,4 $.}

  

%%% Local Variables:
%%% mode: latex
%%% TeX-master: "../am2-19-20"
%%% End:
